\documentclass[11pt,spanish]{article} % Tipo y tamaño de letra del documento.
\usepackage{graphicx}
\usepackage{bm}
\usepackage{hyperref}
\usepackage{natbib}
\graphicspath{{/home/ignacio/Downloads/tarea-1-algoco/template-tarea-1/code/}}

\input{preamble.tex}

\input{author.tex}

\headheight=14pt
\linespread{1.3}
\author{\namestudent}
\pagestyle{fancy}
\fancyhf{}%
\fancyfoot[R]{ \namestudent \\ \rolusm}
\fancyfoot[L]{Campus \campus}
\fancyfoot[C]{\thepage}
\rhead{\sem}
\lhead{INF-221}
\renewcommand{\headrulewidth}{0.4pt}
\renewcommand{\footrulewidth}{0.4pt}
\newbool{programs}
\boolfalse{programs}
\chead{REPORTE TAREA \tnum~}



\title{
  \huge
  \textbf{REPORTE TAREA \tnum~ \\ ALGORITMOS Y COMPLEJIDAD} \\[1ex]
  \emph{\textquote{Más allá de la notación asintótica: Análisis experimental de algoritmos de ordenamiento y multiplicación de matrices.}}
  }


\date{
  \small
  \today\\
  \currenttime
}




\begin{document}
\maketitle
\thispagestyle{fancy}
\vspace{-1.0\baselineskip}

\section{Entorno experimental}
\subsection{Especificaciones del sistema}
\begin{itemize}
  \item \textbf{Sistema operativo:} EndeavourOS x64.
  \item \textbf{Procesador:} Intel i3-N300 3.8GHz.
  \item \textbf{Ram:} 16GB.
  \item \textbf{Compilador:} g++ 15.2.1. 
  \item \textbf{Flags:} -std=c++17, -Wall, -Wextra, -g y -O2
\end{itemize}
\subsection{Datasets}
Los datasets utilizados en este experimento siguen el siguiente formato tanto para el analisis de algoritmos
 de ordenamiento y algoritmos de multiplicación de matrices.
\begin{center}
  \begin{tabular}{|c|c|c|c|c|}
    \hline
    Indice & Tamaño & Tiempo [ns] & Memoria & algoritmo \\
    \hline
    0 & ... & ... & ... & ... \\
    1 & ... & ... & ... & ... \\
    ... & ... & ... & ... & ... \\
    \hline
  \end{tabular}
\end{center}
Fue necesario realizar una limpieza del dataset previo a la generacion de los graficos, ya que el metodo empleado para la 
ejecucion de las pruebas y toma de datos de las mismas genero en el archivo de salida caracteres que al momento de ser 
tratados generaban filas vacias las cuales debian ser removidas.

\section{Resultados y Análisis}
\subsection{Algoritmos de ordenamiento}
\textbf{Introduccion:} Con el fin de poder verificar de forma experimental las complejidades asociadas a algoritmos de 
ordenamiento tales como \textbf{\textit{merge-sort}}, \textbf{\textit{quick-sort}} y 
\textbf{\textit{std::sort}}(libreria estandar de c++), llevamos a cabo mediciones tanto de tiempo de ejecucion como uso 
de memoria de los algoritmos en cuestion para distintos tamaños de entrada ($10^{1}$, $10^{3}$, $10^{5}$ y $10^{7}$), 
para luego poder generar graficos y analizar los resultados con el objetivo de concluir si el comportamiento de los 
algoritmos es el esperado segun sus complejidades teoricas.
\subsubsection{Analisis de tiempo}

\begin{table}[h]
  \small
  \begin{tabular}{|c|c|c|c|c|c|}
    \hline
    & Tamaño entrada & Algoritmo & Tiempo medio de ejecucion & Desviacion estandar & Cantidad de muestras \\
    \hline
    0 & 10 & merge & 6628.888889 & 714.634492 & 18 \\
    1 & 10 & quick & 4066.222222 & 989.871623 & 18 \\
    2 & 10 & std::sort & 4482.611111 & 1134.080044 & 18 \\
    \hline
    3 & 1000 & merge & 378904.944444 & 33217.083191 & 18 \\
    4 & 1000 & quick & 79061.888889 & 19006.587892 & 18 \\
    5 & 1000 & std::sort & 48214.555556 & 17686.456537 & 18 \\
    \hline
    6 & 100000 & merge & 48551985.666667 & 5177993.270268 & 18 \\
    7 & 100000 & quick & 11722067.388889 & 3605245.194886 & 18 \\
    8 & 100000 & std::sort & 6394680.944444 & 2570092.521951 & 18 \\
    \hline
    9 & 10000000 & merge & 6320678299.000000 & 675741581.977279 & 18 \\
    10 & 10000000 & quick & 1594876913.333333 & 420415684.535780 & 18 \\
    11 & 10000000 & std::sort & 881570211.888889 & 347984322.455967 & 18 \\
    \hline
  \end{tabular}
\end{table}
\begin{center}
  \textbf{Tabla 1:} \textit{Tabla comparativa de los datos obtenidos en el experimento.}
\end{center}
Esta tabla muestra el tiempo promedio que le tomo a cada algoritmo terminar su ejecucion para los distintos tamaños de 
entrada.
\clearpage

\begin{figure}[h]
  \centering
  \includegraphics[scale=0.5]{sorting/data/plots/sizeVStime.png}
  \caption{Relacion entre el tamaño de la entrada y tiempo de ejecucion del algoritmo para dicha entrada (escala logaritmica).}
  \label{fig:size vs time}
\end{figure}
El grafico anterior permite hacer una comparacion visual entre algoritmos. Ademas este nos permite de cierta manera 
corroborar las complejidades teoricas determinadas para estos algoritmos de ordenamiento.\newline
En el grafico podemos observar que el comportamiento de estos algoritmos es aproximadamente lineal en escala logaritmica lo cual es consistente con 
la complejidad teorica de estos, la cual para todos es \bm{$O(n\log{n})$} para el caso promedio, por lo que ver 
que los tres algoritmos forman lo mas cercano a una recta en una escala logaritmica es lo esperado.\newline
A pesar de que los tres algoritmos siguen una tendencia casi lineal, podemos notar que las rectas de los algoritmos 
\textbf{\textit{std::sort}} y \textbf{\textit{quick-sort}} son bastante cercanas denotando tiempos de ejecucion promedio 
similares 
para entradas del mismo tamaño, esto debido a que la implementacion de \textbf{\textit{std::sort}} hace uso de \textbf{\textit{introsort}} el cual
 internamente hace uso de \textbf{\textit{quicksort}} junto a otros algoritmos para evitar el peor caso de este ultimo (\bm{$O(n^{2})$}) lo que explicaria la 
 similaridad entre los tiempos promedio de ejecucion para tamaño de entradas menores y una leve superioridad de \textbf{\textit{std::sort}} conforme el tamaño 
 de la entrada aumenta.\newline

Por otro lado en el caso del algoritmo \textbf{\textit{merge-sort}} se midieron tiempos 
promedio de ejecucion mayores, una de la razones de este comportamiento es la naturaleza de estos algoritmos.\newline
Mientras que \textbf{\textit{std::sort}} y \textbf{\textit{quick-sort}} son algoritmos de naturaleza \textit{in-place} lo 
que quiere decir que realizan el ordenamiento sin necesidad de utilizar estructuras auxiliares si no que mas bien logran 
el ordenamiento mediante el uso de recursion, por otro lado \textbf{\textit{merge-sort}} es de naturaleza 
\textit{no in-place} lo que significa que este algoritmo hace uso de estructuras auxiliares, copia de datos a estas mismas, 
mas lectura y escritura de memoria. Los que genera un mayor impacto en el rendimiento.\newline
Esto causa que si bien en la teoria para el caso promedio estos 3 algoritmos tienen la misma complejidad, al momento de 
realizar pruebas reales comienzan a relucir estas diferencias sutiles entre los algoritmos.

\subsubsection{Analisis de memoria}

\begin{table}[h]
  \small
  \begin{tabular}{|c|c|c|c|c|c|}
    \hline
     & Tamaño & Algoritmo & Uso medio de memoria & Desviacion estandar & Cantidad de muestras \\
     \hline
    0 & 10 & merge & 3701.111111 & 60.364470 & 18 \\
    1 & 10 & quick & 3681.777778 & 75.287336 & 18 \\
    2 & 10 & std::sort & 3734.444444 & 83.670377 & 18 \\
    \hline
    3 & 1000 & merge & 3663.555556 & 66.508047 & 18 \\
    4 & 1000 & quick & 3698.666667 & 64.977824 & 18 \\
    5 & 1000 & std::sort & 3697.555556 & 78.947654 & 18 \\
    \hline
    6 & 100000 & merge & 4537.777778 & 67.663141 & 18 \\
    7 & 100000 & quick & 4277.777778 & 45.030128 & 18 \\
    8 & 100000 & std::sort & 4712.888889 & 71.390302 & 18 \\
    \hline
    9 & 10000000 & merge & 82899.777778 & 654.233593 & 18 \\
    10 & 10000000 & quick & 70308.222222 & 548.037735 & 18 \\
    11 & 10000000 & std::sort & 83140.444444 & 573.961899 & 18 \\
    \hline
  \end{tabular}
\end{table}
\begin{center}
  \textbf{Tabla 2:} \textit{Tabla comparativa de los datos obtenidos en el experimento.}
\end{center}
Esta tabla muestra el pico promedio de uso de memoria del algoritmo tomando tambien en cuenta el parseo de las entradas.\newline

\begin{figure}[h]
  \centering
  \includegraphics[scale=0.5]{sorting/data/plots/sizeVSmemory.png}
  \caption{Grafica que relaciona tamaño de entrada con el pico de uso de memoria del algoritmo.}
  \label{fig:memory vs time}
\end{figure}
\clearpage

En el grafico anterior podemos ver como aumenta el pico del uso de memoria conforme aumenta el tamaño de la entrada.\newline
Es importante mencionar que el metodo utilizado para poder medir el uso de memoria no nos permite realizar la medicion
solo del algoritmo si no que mas bien de toda la ejecucion del programa, por lo que en estos picos de memoria tambien se esta
tomando en cuenta la creacion de los arreglos de entrada a partir de archivos de texto.
A pesar de esta causa de ruido en las mediciones, es posible notar como el algoritmo \textbf{\textit{merge-sort}} luego de 
que las entradas son mas grandes que $10^{3}$ el uso de memoria es levemente mayor que el de sus contra partes, esto 
debido a lo mencionado anteriormente sobre la naturaleza \textit{no in-place} del algoritmo, ya que mientras mas grande es la 
entrada mas divisiones se deben hacer sobre el arreglo, lo que implica mas arreglos auxiliares, mas copias en memoria, mas 
lectura de memoria, etc.\newline
Si bien el hecho de medir peaks de memoria para cada ejecucion causa que en el grafico todos los algoritmos utilicen la misma 
cantidad de memoria extra para su funcionamiento si logramos apreciar como \textbf{\textit{merge-sort}} requiere de 
de memoria extra (\bm{$O(n)$}) para llevar a cabo el ordenamiento del arreglo.
\clearpage

\subsection{Algoritmos de multiplicación de matrices}
\textbf{Introduccion:} En esta seccion del experimento realizaremos nuevamente un analisis tanto de tiempo de ejecucion como
de uso de memoria, pero en este caso lo realizaremos sobre algoritmos de multiplicación de matrices.\newline
En concreto analizaremos el comportamiento de 2 algoritmos \textbf{\textit{naive}} y \textbf{\textit{strassen}}, 
siendo el algoritmo naive la opcion mas "\textit{bruta}" para la mutiplicacion ya que consiste en 3 ciclos for anidados, 
mientras que el algoritmo de strassen sigue el enfoque dividir y conquistar lo que en la teoria reduce la complejidad 
temporal del mismo.

\subsubsection{Analisis de tiempo}
\label{tiempo matrices}
\begin{table}[h]
  \small
  \begin{tabular}{|c|c|c|c|c|c|}
    \hline
    & Tamaño & Algoritmo & Tiempo medio de ejecucion & Desviacion estandar & Cantidad de muestras \\
    \hline
    0 & 16 & naive & 184954.777778 & 14339.648402 & 18 \\
    1 & 16 & strassen & 5872281.944444 & 124022.463514 & 18 \\
    \hline
    2 & 64 & naive & 10862325.777778 & 125804.566517 & 18 \\
    3 & 64 & strassen & 302482287.166667 & 8508341.955510 & 18 \\
    \hline
    4 & 256 & naive & 730841801.222222 & 14982636.822267 & 18 \\
    5 & 256 & strassen & 15475774095.277779 & 429441730.023688 & 18 \\
    \hline
    6 & 1024 & naive & 45894626862.722221 & 257400123.342034 & 18 \\
    7 & 1024 & strassen & 760005728144.333374 & 11907145627.682590 & 18 \\
    \hline
  \end{tabular}
\end{table}
\textbf{Tabla 3:} Esta tabla muestra los datos obtenidos de tiempo promedio de ejecucion de los algoritmos para distintos tamaños de entrada.

\begin{figure}[h]
  \centering
  \includegraphics[scale=0.5]{matrix_multiplication/data/plots/sizeVStime.png}
  \caption{Grafica que relaciona tamaño de entrada con tiempo promedio de ejecucion del algoritmo.}
  \label{fig:memory vs time}
\end{figure}
\clearpage

En el grafico anterior podemos apreciar que la tendencia de crecimiento de los algoritmos es polinomial, lo cual es consistente si consideramos las  
  complejidades teoricas para los algoritmos \textbf{\textit{naive}} y \textbf{\textit{strassen}} de multiplicación 
  de matrices tienen complejidades \bm{$O(n^{3})$} y \bm{$O(n^{\log_{2}7}) \approx O(n^{2.81})$} respectivamente.
  Algo interesante que podemos notar es que para todos los tamaños de entrada el algoritmo \textbf{\textit{naive}} 
  tiene tiempos de ejecucion promedio mejores que el algoritmo de \textbf{\textit{strassen}}.\newline
  Esto resulta contradictorio ya que en teoria el algoritmo de \textbf{\textit{strassen}} se estipula que este deberia 
  tener un mejor rendimiento, que a pesar de que la mejora no es demasiada conforme el tamaño de la entrada aumenta 
  deberia quedar mas en evidencia esta mejora pero el grafico muestra una realidad completamente distinta.\newline
  Este fenomeno se puede deber a que el algoritmo de \textbf{\textit{strassen}} en la practica depende fuertemente 
  de la implementacion del mismo, en nuestro caso el uso de estructuras no contiguas tales como \newline
  \verb|std::vector<std::vector<int>>| para representar las matrices y la creacion de multiples submatrices se traducen en un 
  alto nivel de overhead lo que afecta fuertemente al rendimiento del algoritmo, esto sumado al alto costo por copias y 
  asignaciones dinamicas de memoria causan el efecto observado.\newline
  Si bien teoricamente el algoritmo de \textbf{\textit{strassen}} deberia ser mas eficiente que el algoritmo 
  \textbf{\textit{naive}} ocurre que este ultimo para todos los tamaño analizados en promedio entrega la solucion mas 
  rapido.

  \subsubsection{Analisis de memoria}
  \begin{table}[h]
    \small
    \begin{tabular}{|c|c|c|c|c|c|}
      \hline
      & Tamaño & Algoritmo & Uso medio de memoria & Desviacion estandar & Cantidad de muestras \\
      \hline
      0 & 16 & naive & 3692.444444 & 55.185486 & 18 \\
      1 & 16 & strassen & 3670.222222 & 44.398669 & 18 \\
      \hline
      2 & 64 & naive & 3815.777778 & 67.610959 & 18 \\
      3 & 64 & strassen & 3913.111111 & 47.185353 & 18 \\
      \hline
      4 & 256 & naive & 5065.111111 & 49.292511 & 18 \\
      5 & 256 & strassen & 6741.111111 & 50.908094 & 18 \\
      \hline
      6 & 1024 & naive & 24398.222222 & 52.175671 & 18 \\
      7 & 1024 & strassen & 48334.222222 & 94.887894 & 18 \\
      \hline
    \end{tabular}
  \end{table}
\textbf{Tabla 4:} Tabla de datos con el promedio de los picos de uso de memoria de cada algoritmo.\newline
\clearpage
\begin{figure}[h]
  \centering
  \includegraphics[scale=0.5]{matrix_multiplication/data/plots/sizeVSmemory.png}
  \caption{Grafica que relaciona tamaño de entrada con pico de uso de memoria promedio.}
  \label{fig:memory vs time}
\end{figure}
En el grafico anterior podemos notar que el uso de memoria para entradas de tamaño $2^{4}$ y $2^{6}$ para ambos algoritmos es 
practicamente el mismo, esto podria deberse a que en este experimentos se midieron picos de uso de memoria por lo que en este 
grafico igual se esta considerando el uso de memoria asociado a almacenar la entrada hecha mediante texto a un vector 
bidimensional.\newline
Sin embargo podemos notar que para los tamaños de entrada $2^{8}$ y $2^{10}$ el uso de memoria del algoritmo de 
\textbf{\textit{strassen}} crece mas rapido que el uso del algoritmo \textbf{\textit{naive}}, esto debido a lo mencionado 
anteriormente en la seccion \ref{tiempo matrices}.\newline
El algoritmo de \textbf{\textit{strassen}} realiza muchas mas alocaciones dinamicas de memoria y copias en la misma 
para realizar la division en submatrices y multiplicarlas para llegar al resultado final. Esto ademas de tener un 
impacto en el tiempo de ejecucion eleva mucho mas el uso de memoria en comparacion al algoritmo \textbf{\textit{naive}} 
el cual solo hace accesos a memoria mediante indices.
\clearpage

\section{Conclusiones}
En este trabajo se realizó un análisis experimental de algoritmos de ordenamiento y multiplicación de matrices 
con el objetivo de contrastar su comportamiento práctico con sus complejidades teóricas. En el caso de los algoritmos 
de ordenamiento, los resultados obtenidos muestran un crecimiento consistente con una complejidad promedio de 
(\bm$O(n \log n)$), observándose además que \textbf{\textit{std::sort}} y \textbf{\textit{quick-sort}} presentan mejores tiempos de ejecución que \textbf{\textit{merge-sort}}.
 Esta diferencia se explica principalmente por el hecho de que los primeros son algoritmos in-place y presentan un 
menor overhead en memoria, lo que impacta positivamente en el rendimiento. Por otro lado, en la multiplicación de matrices, 
si bien teóricamente el algoritmo de \textbf{\textit{strassen}} posee una mejor complejidad que el algoritmo naive, los resultados 
experimentales muestran que este último es más eficiente para los tamaños de entrada analizados. Esto se debe principalmente 
a factores de implementación, como el uso de estructuras no contiguas y el costo asociado a la creación de submatrices y 
operaciones adicionales, lo que introduce un alto overhead. En conjunto, los resultados evidencian que, aunque la teoría 
permite predecir tendencias generales, el rendimiento real de los algoritmos depende fuertemente de aspectos prácticos como 
la implementación, el uso de memoria y las características del sistema.


\printbibliography

\end{document}
